\section{Decoding tumour phenotype by noninvasive imaging using a quantitative radiomics approach}

Aerts et al. (2014) presented one of the pioneering studies that introduced the concept of radiomics for quantitative tumour characterization based on standard medical imaging \cite{Aerts2014-bg}. The authors extracted 440 handcrafted radiomic features describing tumour intensity, shape, texture, and multiscale wavelet transformations from computed tomography (CT) scans of 1,019 patients with lung and head-and-neck cancer. These features were designed to comprehensively capture tumour phenotype heterogeneity and were subsequently analyzed to evaluate their association with clinical parameters, gene-expression profiles, and patient survival outcomes.

The study employed unsupervised clustering on the Lung1 dataset (422 NSCLC patients) to identify groups of patients with similar radiomic expression patterns. Distinct clusters were found to correlate significantly with tumour stage (T-stage and overall stage) and histology, while no clear association was observed for lymph node or metastasis status. These findings demonstrated that radiomic patterns reflect underlying phenotypic and clinical characteristics of tumours.

To assess the prognostic value of radiomic features, the authors performed survival analysis using Kaplan-Meier curves across independent datasets (Lung1, Lung2, H\&N1, and H\&N2). Using a median feature threshold derived from the training set, they validated performance across cohorts without retraining. More than half of the radiomic features (54\%) were significantly associated with overall survival in the Lung2 validation set, and 15\% retained prognostic significance across all three external datasets. Features describing intra-tumour heterogeneity were strongly linked to poorer survival outcomes, while shape features indicating more compact tumours correlated with improved prognosis.

Aerts et al. further evaluated the reproducibility of radiomic features using the publicly available RIDER dataset (test-retest CT scans) and a multiple-delineation dataset annotated by five oncologists. This analysis demonstrated that a substantial subset of features was stable across repeated acquisitions and manual segmentations. The final prognostic radiomic signature thus represented a set of robust and reproducible imaging biomarkers associated with survival and gene-expression patterns. Overall, this landmark study established radiomics as a quantitative, non-invasive bridge between imaging and tumour biology, providing a foundation for future radiogenomic and AI-based approaches in precision oncology.
